% SPDX-FileCopyrightText: Copyright (c) 2024-2026 Yegor Bugayenko
% SPDX-License-Identifier: MIT

\section{Texts}\label{sec:strings}

All objects presented in this Section belong to the \ff{tt} package.

The \adeff{sscanf} object encapsulates two \deff{string} objects:
  a format and a content.
It behaves like a tuple of scanned data objects.
For example, this code parses a hexadecimal number from the console:

\begin{ffcode}
at.
  tt.sscanf
    "%x"
    io.stdin
  0
\end{ffcode}

The \adeff{sprintf} object builds a string according to the format
  provided and compliant with \citet[Chapter 5]{posix}, for example:

\begin{ffcode}
tt.sprintf *1
  "Hi, %s! The weight is %0.2f."
  "Jeff"
  3.14
\end{ffcode}

A few more decorators of \deff{string}:

\begin{itemize}
    \item \deff{trimmed}: removes both leading and trailing spaces
    \item \deff{trimmed-left}: removes leading spaces
    \item \deff{trimmed-right}: removes trailing spaces
    \item \deff{joined}: joins a tuple of strings with this one as a glue
    \item \deff{chained}: concatenates itself with other strings
    \item \deff{repeated}: repeats itself a given number of times
    \item \deff{slice}: takes a piece of itself as another string
    \item \deff{at}: retrieves the symbol at a given index
    \item \deff{contains}: checks whether it contains a substring
    \item \deff{contains-all}: checks whether it contains all given substrings
    \item \deff{contains-any}: checks whether it contains any given substring
    \item \deff{starts-with}: checks whether it starts with a substring
    \item \deff{ends-with}: checks whether it ends with a substring
    \item \deff{index-of}: finds the first occurrence of a substring
    \item \deff{last-index-of}: finds the last occurrence of a substring
    \item \deff{up-cased}: makes it upper case
    \item \deff{low-cased}: makes it lower case
    \item \deff{replaced}: finds and replaces a substring
    \item \deff{split}: breaks it into a tuple of strings
    \item \deff{nsplit}: breaks it into at most a given number of strings
    \item \deff{as-number}: converts itself to \deff{number}
    \item \deff{as-ascii}: reads a single character as its ASCII code
    \item \deff{is-alpha}: checks whether all characters are letters
    \item \deff{is-digit}: checks whether all characters are digits
    \item \deff{is-alphanumeric}: checks whether all characters are letters or digits
    \item \deff{is-ascii}: checks whether all characters are ASCII
\end{itemize}

The \deff{string-buffer} object allows building a string incrementally,
  appending pieces through its \deff{with} attribute.
For example, this code builds the string \ff{"Hello, Jeff"}
  and afterwards behaves as a plain \deff{string}:

\begin{ffcode}
tt.string-buffer "Hello"
  .with ", "
  .with "Jeff"
\end{ffcode}

The \deff{regex} object is an abstraction of a regular expression
  in Perl-compatible format, with full support of Unicode.
The \deff{matches} attribute is \deff{true} if the given text
  matches the pattern.
The \deff{match} attribute is a matcher over the given text: its
  \deff{next} attribute is the first matched block, each block exposes
  \deff{text}, \deff{start}, and \deff{group}, and its own \deff{next}
  attribute is the following block.

\begin{ffcode}
tt.regex "/[a-z]+/" > r
r.match "!hello!world!" > m
eq.
  m.next.text
  "hello"
\end{ffcode}
